\documentclass[11pt]{article}
\usepackage[top=2.54cm, bottom=2.54cm, left=2.75cm, right=2.75cm]{geometry}
\linespread{1.2}
\setlength{\parindent}{0cm}
\usepackage{amsmath}
\usepackage{amssymb}
\usepackage{amsthm}
\usepackage{datetime}
\usepackage{graphicx}
\usepackage{float}
\usepackage{tikz}
\usepackage{pgfplots}
\usepackage{hyperref}
\usepackage{subcaption}

\bibliographystyle{ieeetr}
\usepackage{bibentry}
\nobibliography*

\newcommand{\vd}{\mathrm{d}} %Vertical d (for derivatives)

\begin{document}
	\pagenumbering{arabic}
	
	\begin{titlepage}
		\begin{center}
			{\Huge Monte Carlo Simulations of The Ising Model}\\[0.5cm]
			\textit{Haodong Chang and Tavis O'Reilly}\\[0.3cm]
			\textit{11505303 and 10903943}\\[0.3cm]
			Department of Physics and Astronomy\\[0.3cm]
			University of Manchester\\[0.3cm]
			PHYS20872 Theory Computing Project\\[0.3cm]
			\shortmonthname[\the\month]  \the\year \\[4cm]
			
		\end{center}
		
		{\Large \textbf{Abstract}}\\[0.3cm]
		The goal of this project was to use a Metropolis-Hastings algorithm, a type of Markov chain Monte Carlo method, to simulate the Ising model for ferromagnetism. In two dimensions, an exact analytic solution exists\cite{onsager_solution} that shows a phase transition at some critical temperature $T_c$. We wrote a computer simulation that implements a Metropolis-Hastings algorithm to find the average energy, heat capacity and magnetisation as a function of temperature. Comparing our results in two dimensions to the analytic solution, we found them to closely match. We then extended our simulation to cases for which there is no known exact solution such as the two dimensional Ising model with an external magnetic field and the three dimensional case.
		
		\vfill
		[1] \bibentry{onsager_solution}
	\end{titlepage}
	
	\pagenumbering{gobble}
	\clearpage
	\pagenumbering{arabic}
	\setcounter{page}{2}
	
	\newpage
	
	\section{Introduction}
	
	The Ising model is a simple model of ferromagnetism, consisting of a lattice of atomic spins that can either be "up" or "down". These spins only interact with their nearest neighbours, and with any external magnetic field. The system "prefers" to be in its lowest energy state (all spins aligned), however this average energy will increase with temperature due to the random energy transfers from the surroundings. There exists an exact analytic solution to the model in two dimensions (for zero magnetic field)\cite{onsager_solution}, found by Lars Onsager in 1944; in three dimensions there have been many attempts at a solution\cite{das1970exact, lou2000three, zhang2007conjectures, zhang2021exact} but each has been shown to have errors. Instead, we used numerical methods to solve the model for the temperature dependence of energy, magnetisation and heat capacity. This allows the model to be trivially extended to solve for cases such as an external magnetic field and in three dimensions.
	
	\section{Theory}
	
	For a set of $N$ spins $\left\{s_i\right\}_{i=1}^N$ arranged in a lattice, where the spin is $+1$ for "up" and $-1$ for "down", the Hamiltonian (or energy)	of the system $H$ is given by\cite{Baierlein_1999} % use $H$ to avoid confusion with the average energy $E$
	\begin{equation}
		H\left(\left\{s_i\right\}\right) = -J\sum_{\left<i,j\right>}s_is_j - h\sum_i s_i,
	\end{equation}
	where $J$ is the interaction strength between adjacent spins, $h$ is the strength of the external magnetic field and $\sum_{\left<i,j\right>}$ indicates the sum over all nearest neighbours $j$ of the $i$th spin. Some simplifications have been made as in principle $J$ could not be constant. \\
	
	We will be dealing with the canonical ensemble, where the temperature is fixed. The key to solving the model is its partition function $Z$ given by
	\begin{equation} \label{eq:PartitionFunction}
		Z = \sum_{\{s_i\}} e^{-\beta H(\{s_i\})},
	\end{equation}
	where $\sum_{\{s_i\}}$ indicates the sum over all possible sets of $N$ spins and $\beta = \frac{1}{k_B T}$, where $k_B$ is the Boltzmann constant and $T$ is the temperature. The canonical ensemble assigns a probability to each possible state $\{s_i\}$ of
	\begin{equation} \label{eq:Probability}
		P(\{s_i\}) = \frac{e^{-\beta H(\{s_i\})}}{Z}
	\end{equation}

	\subsection{One Dimension}
	
	In one dimension, the model can be solved in relatively few steps. In this case, the Hamiltonian of the system for $N$ spins in a line becomes
	\begin{equation}
		H(\{s_i\}) = -J \sum_{i=1}^{N-1} s_i s_{i+1} -h\sum_{i=1}^N s_i.
	\end{equation}
	Alternatively, periodic boundary conditions can be applied so that $s_{N+1} = s_1$. This is advantageous as every spin then has the same behaviour, and in the limit of $N \to \infty$ this is equivalent. Our discussions and simulations below will be based on periodic boundary conditions. In this case, the energy of the system becomes
	\begin{equation} \label{eq:Energy1D}
		H(\{s_i\}) = -J \sum_{i=1}^{N} s_i s_{i+1} -h\sum_{i=1}^N s_i.
	\end{equation}

	Before dealing with such a relatively complicated system, let's consider a simple Ising model with only two spins. The Hamiltonian is
	\begin{equation} \label{eq:Hamiltonian1DGeneral2Spins}
		H(s_1, s_2) = -J s_1 s_2 - J s_2 s_1 - hs_1 - hs_2
	\end{equation}

	Define a matrix $P$, we can see why we need it later.
	\begin{equation} \label{eq:MatrixP}
		P = \begin{pmatrix}
			e^{\beta(J+h)} & e^{-\beta J} \\
			e^{-\beta J} & e^{\beta(J-h)}
		\end{pmatrix}
	\end{equation}

	Then the partition function is
	\begin{equation} \label{eq:PartitionFunction1DGeneral2Spins}
		Z = \sum_{s_1, s_2} e^{-\beta H(s_1, s_2)} = \mathrm{Tr}(P \cdot P)
	\end{equation}
	Here $\mathrm{Tr}$ denotes the trace of a matrix.

	One can easily verify the above result by expanding the matrix product, but may be curious about why we need such a matrix $P$. Actually, the matrix $P$ indicates the "connection" between the two spins.

	Suppose that the first row and column of $P$ represent a state $s = +1$, while the second row and column represent another state $s = -1$. We can write down the elements of $P$ as
	\begin{equation} \label{eq:MatrixPElements}
		P_{+,+} = e^{\beta(J+h)}, \quad P_{+,-} = P_{-,+} = e^{-\beta J}, \quad P_{-,-} = e^{\beta(J-h)}
	\end{equation}

	The contribution of the connection between two spins to the partition function is
	\begin{equation} \label{eq:PartitionFunction1DGeneral2SpinsConnection}
		P_{s_1, s_2} = e^{\beta(J s_1 s_2 + \frac12h(s_1 + s_2))}
	\end{equation}
	The factor $\frac12$ here is because each spin will be counted twice in two connections it has.

	We can write down the partition function of such a system with $2$ spins as
	\begin{equation} \label{eq:PartitionFunction1DGeneral2SpinsConnectionSum}
		\begin{aligned}
		Z &= \sum_{s_1, s_2} e^{-\beta H(s_1, s_2)} = \sum_{s_1, s_2} e^{\beta(2J s_1 s_2 + h(s_1 + s_2))} \\
		&= \sum_{s_1, s_2} e^{\beta(J s_1 s_2 + \frac12h(s_1 + s_2))} \cdot e^{\beta(J s_1 s_2 + \frac12h(s_1 + s_2))} \\
		&= \sum_{s_1, s_2} P_{s_1, s_2} \cdot P_{s_2, s_1} \\
		&= \mathrm{Tr}(P \cdot P)
		\end{aligned}
	\end{equation}

	It's easy to generalize the result to a system with $N$ spins.
	\begin{equation} \label{eq:PartitionFunction1DGeneral}
		\begin{aligned}
			Z &= \sum_{\{s_i\}} e^{-\beta H(\{s_i\})} = \sum_{\{s_i\}} \prod_{i=1}^{N} e^{\beta(J s_i s_{i+1} + \frac12h(s_i + s_{i+1}))} \\
			&= \sum_{\{s_i\}} \prod_{i=1}^{N} P_{s_i, s_{i+1}} = \mathrm{Tr}(P^N) \\
		\end{aligned}
	\end{equation}

	Now our task is simplified into calculating the trace of the $N$-th power of a $2 \times 2$ matrix $P$. Mathematicians have already solved this problem by diagonalizing the matrix
	\begin{equation} \label{eq:MatrixPDiagonal}
		P = U \cdot D \cdot U^\dagger
	\end{equation}
	Here $U$ is a unitary matrix, which means $U^\dagger U = I$, and $D$ is a diagonal matrix. It has been proven that such diagonalization is always possible for a Hermitian matrix like $P$.

	We only care about the trace of $P^N$
	\begin{equation} \label{eq:MatrixPNTrace}
		\mathrm{Tr}(P^N) = \mathrm{Tr}(U \cdot D^N \cdot U^\dagger) = \mathrm{Tr}(D^N)
	\end{equation}

	For a diagonal matrix $D = \mathrm{diag}(\lambda_+, \lambda_-)$, its $N$-th power is also a diagonal matrix $D^N = \mathrm{diag}(\lambda_+^N, \lambda_-^N)$. Thus, we have
	\begin{equation} \label{eq:MatrixPNTraceDiagonal}
		\mathrm{Tr}(P^N) = \lambda_+^N + \lambda_-^N
	\end{equation}

	Now we only need to find $\lambda_\pm$, which are the eigenvalues of $P$. I used Mathematica to calculate them, and the result is
	\begin{equation} \label{eq:MatrixPEigenvalues}
		\lambda_\pm = e^{\beta J} \left[\cosh(\beta h) \pm \sqrt{e^{-4\beta J}+\sinh^2(\beta h)}\right]
	\end{equation}

	Finally, we have the partition function of a 1D Ising model with $N$ spins
	\begin{equation} \label{eq:PartitionFunction1DGeneralFinal}
		\begin{aligned}
			Z &= \mathrm{Tr}(P^N) = \lambda_+^N + \lambda_-^N \\
			&= e^{N \beta J} \left\{\left[\cosh(\beta h) + \sqrt{e^{-4\beta J}+\sinh^2(\beta h)}\right]^N
			 + \left[\cosh(\beta h) - \sqrt{e^{-4\beta J}+\sinh^2(\beta h)}\right]^N\right\}
		\end{aligned}
	\end{equation}

	In the limit $N \to \infty$, we have $\lambda_+^N \gg \lambda_-^N$, and the partition function becomes
		\begin{equation} \label{eq:PartitionFunction1DGeneralLimit}
			Z \approx \lambda_+^N = e^{N \beta J} \left[\cosh(\beta h) + \sqrt{e^{-4\beta J}+\sinh^2(\beta h)}\right]^N
		\end{equation}
	It will be much easier for us to calculate energy $E$, heat capacity $C_V$, magnetisation $M$, and other thermodynamic properties under such approximation.

	\subsection{2D Ising Model}
	
	Onsager solved the 2D Ising model without external magnetic field ($h = 0$) in 1944\cite{onsager_solution}, but solution for $h \neq 0$ is still yet to be found. He used transfer matrix method, like what we have done in the 1D case, but much more complicated. Several years later, a combinatorial method was developed by Kac and Ward\cite{KacWard1952}, which is much simpler than the transfer matrix method, but still involves great details. We will only present a method to calculate the critical temperature $T_{cr}$ due to Kramers and Wannier\cite{KramersWannier1941}, and then we will give the partition function and energy without proof.

	By definition, the partition function of a 2D Ising model without external field is
	\begin{equation} \label{eq:PartitionFunction2D}
		Z = \sum_{\{s_i\}} \prod_{\langle i,j \rangle} e^{\beta J s_i s_j}
	\end{equation}
	It's easy to see that if the sign of $J$ is changed, the partition function remains the same. Thus, we can set $J > 0$ without loss of generality.
	% Not quite easy to see. Consider flipping every other spin, but this requires $N$ to be even under PBC. Maybe it doesn't matter when $N$ is big enough.

	Since $s_i s_j = \pm 1$, we have
	\begin{equation} \label{eq:PartitionFunction2DTransformingTrick}
		\begin{aligned}
			e^{\beta J s_i s_j} &= \cosh(\beta J) + s_i s_j \sinh(\beta J) \\
			&= \cosh(\beta J)(1 + s_i s_j \tanh(\beta J))
		\end{aligned}
	\end{equation}

	There are $2N$ pairs of nearest neighbours $\langle i,j \rangle$ in a 2D lattice with $N$ spins. Let $t = \tanh(\beta J)$, and let's try to expand the product in equation \eqref{eq:PartitionFunction2D}. Take each term after expansion as a graph $G = (V, E)$ (so there are $2^{2N}$ graphs in total), where $V$ is the set of spins, and $E$ is the set of $\langle i,j \rangle$ pairs which contribute $t s_i s_j$ to the product. Then we have
	\begin{equation} \label{eq:PartitionFunction2DTransformed}
		\begin{aligned}
			Z &= \sum_{\{s_i\}} \prod_{\langle i,j \rangle} \cosh(\beta J)(1 + t s_i s_j) \\
			&= \cosh^{2N}(\beta J)\ \sum_{\{s_i\}}\ \sum_{G\ \text{with}\ |E|=L} t^L \prod_{\langle i,j \rangle \in E} s_i s_j \\
			&= \cosh^{2N}(\beta J) \sum_{G\ \text{with}\ |E|=L} t^L \sum_{\{s_i\}} \prod_{\langle i,j \rangle \in E} s_i s_j \\
			&= 2^N \cosh^{2N}(\beta J) \sum_{L} t^L f(L)
		\end{aligned}
	\end{equation}
	where $f(L)$ is the number of "closed graphs" with $L$ edges. Here we define a closed graph as a graph with an even number of edges being incident to each vertex, which means that the sum $\sum_{\{s_i\}} \prod_{\langle i,j \rangle \in E} s_i s_j$ is not $0$ but $2^N$.

	Now our task has been simplified into calculating $f(L)$. Given a closed graph on a 2D lattice, we can always define "inside" and "outside" of the graph. Label the midpoints of the lattice to be $x_i=+1$ if they are inside the graph, and $x_i=-1$ if outside. Then link each pair of nearest neighbours $x_i=+1$ and $x_j=-1$ by an edge, and we will get a new graph $G'$ with $L$ edges. For each edge in $G$, there is one and only one edge in $G'$ intersecting it.

	% (To be added: a figure)

	If we take the new graph $G'$ as a lattice with a set of spins $\{x_i\}$, strength of interaction $J$, and $\beta'$ representing the temperature, then the partition function of such a lattice defined by $G'$ is
	\begin{equation} \label{eq:PartitionFunction2DNewLattice}
		\begin{aligned}
			Z' &= \sum_{\{x_i\}} \prod_{\langle i,j \rangle} e^{\beta' J x_i x_j} \\
			&= \sum_{\{x_i\}} e^{-L\beta' J} \cdot e^{(2N-L)\beta' J} g(L) \\
			&= e^{2N\beta' J} \sum_{L} e^{-2L\beta' J} g(L)
		\end{aligned}
	\end{equation}
	where $g(L)$ is the number of ways to arrange $L$ edges in graph $G'$, which means the number of ways to arrange $L$ pairs of nearest neighbours $x_i=+1$ and $x_j=-1$.
	
	Considering the correspondence between $G$ and $G'$, we have $f(L) = g(L)$. Let $e^{-2\beta' J} = t$, remembering that $t = \tanh(\beta J)$. Join equations \eqref{eq:PartitionFunction2DTransformed} and \eqref{eq:PartitionFunction2DNewLattice}
	\begin{equation} \label{eq:PartitionFunction2DComparison}
		Z = 2^N \cosh^{2N}(\beta J) \tanh^{-N}(\beta J) Z'
	\end{equation}

	Now let's take $Z$ as a function of $\beta$, and $Z'$ as a function of $\beta'$. When $\beta$ is finite, $\beta' J = -\frac12 \ln{\tanh(\beta J)}$ is also finite. From equation \eqref{eq:PartitionFunction2DComparison}, we can see that if $Z(\beta)$ has a singularity at $\beta = \beta_{cr}$, then $Z'(\beta')$ has a singularity at $\beta' = \beta'_{cr}$.
	If we assume that there is only one critical temperature $T_{cr}$, which means only one singularity, we will get
	\begin{equation} \label{eq:Singularity}
		\beta_{cr} = \beta'_{cr}
	\end{equation}

	Solution to equation \eqref{eq:CriticalTemperature} yeilds the critical temperature
	\begin{equation} \label{eq:CriticalTemperature}
		k_B T_{cr} = \frac{2}{\ln(1+\sqrt{2})}J \approx 2.269J
	\end{equation}

	We have found $T_{cr}$ of the 2D Ising model, but if we want to calculate the partition function, we still need to figure out $f(L)$, which is a combinatorial problem. You can refer to Feynman's book\cite{feynman1972statistical} for more details.

	The results we are listing below are obtained from the Onsager solution\cite{onsager_solution}. Let $\lambda$ be the maximum eigenvalue of transfer matrix $P$, like in the 1D case. Then when $N$ is large, the partition function is

	\begin{equation}
		\begin{split}
			Z &= \lambda^N \\
			\ln{\lambda} &= \ln{\left(2 \cosh{2\beta J}\right)} + \frac{1}{\pi}\int_0^{\frac{\pi}{2}} \ln{\left[\frac{1}{2}\left(1 + \sqrt{1 - K^2\sin^2{w}}\right)\right]\;\vd w} \\
			K &= \frac{2\sinh{2\beta J}}{\left(\cosh{2\beta J}\right)^2}
		\end{split}
	\end{equation}
	
	The average energy is
	\begin{equation}
		\begin{split}
			E &= - \frac{\partial \ln{Z}}{\partial \beta} = -N\frac{\partial \ln{\lambda}}{\partial \beta}
		\end{split}
	\end{equation}

	\subsection{Metropolis-Hastings Algorithm}
	\label{metropolis_algorithm}
	
	The Metropolis-Hastings algorithm is a Markov Chain Monte Carlo method used to sample from a probability distribution. In this case of the Ising model, the algorithm is as follows:
	\begin{enumerate}
		\item Choose a spin at random from the grid.
		\item Calculate the energy change $\Delta H$ if this spin were to flip.
		\item If $\Delta H \leq 0$ flip the spin, otherwise flip the spin with probability $e^{-\beta\Delta H}$.
		\item Repeat this process.
	\end{enumerate}

	To validate the algorithm, we shall prove that when the system reaches equilibrium, we will have a canonical ensemble, i.e. a probability distribution given by equation \eqref{eq:Probability}.
	
	Suppose we have an Ising model with $N$ spins. Let $A, B$ be two states of the system. The Markov Chain is described by a transition matrix $\pi$ with elements $\pi_{AB}$, which is the probability of transitioning from state $A$ to state $B$. Only when there is no more than one spin flipped between states $A$ and $B$ can $\pi_{AB}$ be non-zero. For such $A \neq B$, we have
	\begin{equation}
		\pi_{AB} = \begin{cases}
			\frac{1}{N} & H(A) \geq H(B) \\
			\frac{1}{N} e^{-\beta(H(B) - H(A))} & H(A) < H(B)
		\end{cases}
	\end{equation}
	and $\pi_{AA} = 1 - \sum_{B \neq A} \pi_{AB}$.

	When the equilibrium is reached, for any two states $A, B$, we can expect a balance of probability fluxes
	\begin{equation}
		P(A) \pi_{AB} = P(B) \pi_{BA}
	\end{equation}
	By simplifying the above equation we get
	\begin{equation}
		\frac{P(A)}{P(B)} = \frac{\pi_{BA}}{\pi_{AB}} = \frac{e^{-\beta H(A)}}{e^{-\beta H(B)}}
	\end{equation}
	Introduce $Z$ for normalization
	\begin{equation}
		\begin{aligned}
			P(A) &= \frac{e^{-\beta H(A)}}{Z} \\
			\text{where}\ Z &= \sum_A e^{-\beta H(A)}
		\end{aligned}
	\end{equation}
	which is exactly the canonical ensemble.

	\section{Implementation}
	
	The first step was to create a grid of spins. We chose to have the spins randomly distributed but with a bias towards the "up" direction. This ensured the magnetisation, which is sensitive to initial conditions, approached the same average each time.\\
	
	The next step is then to repeatedly apply the algorithm described in section \ref{metropolis_algorithm}. Starting in two dimensions, we ran some initial tests of the model to observe the behaviour of the spins as the algorithm iterates, as shown in figure \ref{fig:ising_grid}. This nicely demonstrates the tendency for groups of aligned spins to form.
	\begin{figure}[H]
		\begin{center}
			\includegraphics[width=0.5\textwidth]{./img/ising-simulation.png}
		\end{center}
		\caption{A two-dimensional Ising Model simulation on a $500\times 500$ grid with parameters $\beta = 1.0$ and $h = 0$. The two colours correspond to the two different values of spin.}
		\label{fig:ising_grid}
	\end{figure}
	
	In order to find the temperature dependence of the energy, magnetisation and heat capacity, we needed to run the same simulation for a number of different values of $\beta$. Each simulation needs to be run for a sufficient number of iterations such that its mean values reach equilibrium. The energy and magnetisation will continue to change, but the average should approach a specific value. To test for this, we applied the following stopping condition:
	\begin{itemize}
		\item For each block of $n$ iterations:
		\begin{enumerate}
			\item Calculate the mean energy and magnetisation values for the current block.
			\item Compare the change of these mean values from the current block to those of the previous block.
			\item If both values have changed by less than three times the standard deviation of the mean $\displaystyle\left(\frac{\sigma}{\sqrt{n}}\right)$, it is considered to be at equilibrium and the simulation stops.
		\end{enumerate}
		\item There is also a set maximum iterations $N_\text{max}$ to ensure the algorithm doesn't continue ad infinitum.
	\end{itemize}
	It is important to ensure the block size $n$ is large enough that the probability of the stopping condition triggering erroneously is minimised. \\
	
	Once the simulation has finished for the range of temperature values, the calculated energy and magnetisation along with their respective temperatures are recorded in a CSV file. From the energy data, the heat capacity values can be calculated using $C_V = \frac{\partial E}{\partial T}$. To calculate this gradient, a least-square fit was used along small intervals of energy values in order to estimate the gradient. This reduces errors from random variations in the calculated energy values.
	
	\section{Results}
	
	First, we ran a 2D simulation with grid size $1000 \times 1000$ and no external magnetic field, obtaining the results shown in figure \ref{fig:2d_1000_result}.
	\begin{figure}[H]
		\begin{subfigure}{\textwidth}
			\begin{center}
				\begin{tikzpicture}
					\begin{axis}[width=0.8\textwidth, height=0.5\textwidth, xlabel=$k_BT$, ylabel=$E/N$]
						\addplot [mark size=0pt,blue] table[
						col sep=comma,
						x=temp,
						y=energy]
						{csv/onsager-energy.csv};
						
						\addplot [mark=*,mark size=0.5pt,black] table[
						col sep=comma,
						x=temp,
						y=energy,
						only marks]
						{csv/energy-1000.csv};
					\end{axis}
				\end{tikzpicture}
			\end{center}
			\caption{A plot showing average energy per particle against $k_BT$.}
		\end{subfigure}
	
		\begin{subfigure}{0.45\textwidth}
			\begin{center}
				\begin{tikzpicture}
					\begin{axis}[width=\textwidth, height=0.8\textwidth, xlabel=$k_BT$, ylabel=$M/N$]
						\addplot [mark size=0pt,blue] table[
						col sep=comma,
						x=temp,
						y=magnetisation]
						{csv/onsager-magnetisation.csv};
						
						\addplot [mark=*,mark size=0.5pt,black] table[
						col sep=comma,
						x=temp,
						y=magnetisation,
						only marks]
						{csv/magnetisation-1000.csv};
					\end{axis}
				\end{tikzpicture}
			\end{center}
			\caption{A plot showing average magnetisation per particle against $k_BT$.}
		\end{subfigure}
		\begin{subfigure}{0.45\textwidth}
			\begin{center}
				\begin{tikzpicture}
					\begin{axis}[width=\textwidth, height=0.8\textwidth, xlabel=$k_BT$, ylabel=$C_V/N$]
						\addplot [mark size=0pt,blue] table[
						col sep=comma,
						x=temp,
						y=heat-cap]
						{csv/onsager-heat-cap.csv};
						
						\addplot [mark=*,mark size=0.5pt,black] table[
						col sep=comma,
						x=temp,
						y=heat-cap,
						only marks]
						{csv/heat-cap-1000.csv};
					\end{axis}
				\end{tikzpicture}
			\end{center}
			\caption{A plot showing average heat capacity per particle against $k_BT$.}
		\end{subfigure}
	
		\caption{Plots of data from a 2D Ising model simulation with grid size $1000 \times 1000$ and no external magnetic field. In each plot, the blue line is the expected result using Onsager's solution, and the black points are the data calculated by our simulation.}
		\label{fig:2d_1000_result}
	\end{figure}
	
	From these results we can see that they closely follow the exact solution, indicating that using a Metropolis-Hastings algorithm is a valid method for predicting the behaviours of the model. The results for heat capacity don't have the same sharp peak as the exact solution. Since this peak tends to infinity in the exact solution, our value is limited by both the resolution of temperature values we were using in our simulation and the size of the grid.
	
	- Different magnetic fields
	- 3D case?
	
	\section{Conclusion}
	
	Conclude results.
	
	\bibliography{references} 

\end{document}